%! Polar Function Graphs — Unit 3, Lesson 3.14 — Warm-Up
\documentclass[10pt]{article}
\usepackage{precalculus-article}
\usepackage{precalculus-boxes}
\newcommand{\TallMath}[1]{$\displaystyle #1\rule[-1.4em]{0pt}{3.2em}$}

\begin{document}

\pageheader{Unit 3, Lesson 3.14}{Warm-Up}
\namedateperiod

\vspace{0.10in}

\begin{notesbox}{Spiral Review}

\textbf{1.}\quad Evaluate each expression. Use your knowledge of the unit circle.

\vspace{4pt}
\renewcommand{\arraystretch}{1.6}
\begin{tabular}{@{}p{0.22\linewidth}p{0.22\linewidth}p{0.22\linewidth}p{0.22\linewidth}@{}}
(a)\;$\sin(0) = $ \blank{1.2cm}
& (b)\;$\sin\!\left(\dfrac{\pi}{2}\right) = $ \blank{1.2cm}
& (c)\;$\sin(\pi) = $ \blank{1.2cm}
& (d)\;$\sin\!\left(\dfrac{3\pi}{2}\right) = $ \blank{1.2cm}
\end{tabular}

\vspace{0.16in}
\textbf{2.}\quad The rectangular graph of $r = 2\sin\theta$ is shown below.
Use it to answer the questions.

\vspace{4pt}
\begin{center}
\begin{tikzpicture}[x=1.1cm, y=0.7cm]
  \draw[linegray, very thin] (-0.2,-2.4) grid[xstep=0.5, ystep=0.5] (6.7,2.4);
  \draw[->] (-0.2,0) -- (6.9,0) node[right]{\small $\theta$};
  \draw[->] (0,-2.5) -- (0,2.6) node[above]{\small $r$};
  \foreach \x/\lbl in {1.571/$\frac{\pi}{2}$, 3.142/$\pi$, 4.712/$\frac{3\pi}{2}$, 6.283/$2\pi$} {
    \draw (\x,0.10) -- (\x,-0.10);
    \node[below, font=\scriptsize] at (\x,-0.10) {\lbl};
  }
  \foreach \y/\lbl in {-2/{$-2$}, -1/{$-1$}, 1/{$1$}, 2/{$2$}} {
    \draw (0.07,\y) -- (-0.07,\y);
    \node[left, font=\scriptsize] at (-0.07,\y) {\lbl};
  }
  \draw[navy, thick, domain=0:6.3, samples=100]
    plot (\x, {2*sin(deg(\x))});
  \node[right, font=\small, navy] at (6.4, 0.4) {$r = 2\sin\theta$};
\end{tikzpicture}
\end{center}

\vspace{2pt}
(a)\quad At $\theta = \dfrac{\pi}{2}$, the value of $r$ is \blank{1.5cm}.

\vspace{6pt}
(b)\quad At $\theta = \pi$, the value of $r$ is \blank{1.5cm}.

\vspace{6pt}
(c)\quad On the interval $\theta \in (0, \pi)$, is $r$ positive or negative? \blank{2.5cm}

\vspace{0.14in}
\textbf{3.}\quad In polar coordinates, a point with $r = 0$ is located at the
\underline{\hspace{3cm}} (origin / pole). If the graph of $r = f(\theta)$ passes
through the pole, then $f(\theta) = $ \blank{1cm} at that angle.

\vspace{0.14in}
\textbf{4.}\quad A polar function assigns a \underline{\hspace{3cm}} (output) to each
\underline{\hspace{3cm}} measure (input) $\theta$.

\end{notesbox}

\end{document}
